\documentclass[11pt]{article}

\usepackage[a4paper,margin=1in]{geometry}
\usepackage[T1]{fontenc}
\usepackage{lmodern}
\usepackage{amsmath,amssymb,booktabs,array,enumitem,hyperref}

\setlength{\parindent}{0pt}
\setlength{\parskip}{0.7em}

\newcommand{\T}{\mathbf{T}}
\newcommand{\R}{\mathbf{R}}
\newcommand{\p}{\mathbf{p}}
\newcommand{\q}{\mathbf{q}}
\newcommand{\J}{\mathbf{J}}
\newcommand{\zhat}{\hat{\mathbf{z}}}
\newcommand{\Rotz}{\operatorname{Rot}_z}
\newcommand{\Rotx}{\operatorname{Rot}_x}
\newcommand{\Transz}{\operatorname{Trans}_z}
\newcommand{\Transx}{\operatorname{Trans}_x}

\title{SO101 3-DOF Arm: Simulation Model and Teaching Kinematics}
\author{A self-contained teaching note}
\date{}

\begin{document}
\maketitle

\section{Purpose}
This document explains the current 3-degree-of-freedom arm project from first
principles. It is written for a reader who has not seen the code before.

The project intentionally has two kinematic layers:
\begin{enumerate}[leftmargin=1.3em]
  \item An intact MuJoCo simulation model. This preserves the recognizable
        SO101 visual structure and uses the STL mesh assets.
  \item A clean Denavit-Hartenberg teaching model. This gives a compact
        parameter table, forward kinematics, Jacobian, and inverse kinematics.
\end{enumerate}

The MuJoCo model is the visual and simulation model. The DH model is the
mathematical teaching model. They serve different purposes and should not be
forced into a single model by moving mesh geometry into DH frames.

\section{Abbreviations}
\begin{center}
\begin{tabular}{@{}ll@{}}
\toprule
Abbreviation & Meaning \\
\midrule
3-DOF & Three degrees of freedom, meaning three independent joint variables \\
SO101 & The original arm family from which the geometry/assets were derived \\
MJCF & MuJoCo XML format, the model-description format used by MuJoCo \\
MuJoCo & Multi-Joint dynamics with Contact, the physics simulator \\
STL & Stereolithography mesh file, used here for visual robot parts \\
TCP & Tool center point; here this is the \texttt{gripperframe} site used in FK/IK \\
EE & End effector, the final fixed tool and gripper subassembly \\
DH & Denavit-Hartenberg, a standard robot-link parameter convention \\
FK & Forward kinematics, computing pose from joint angles \\
IK & Inverse kinematics, computing joint angles from a desired pose/position \\
COM & Center of mass \\
\bottomrule
\end{tabular}
\end{center}

\section{Coordinate convention}
The simulator uses a right-handed world coordinate system. In the project:
\begin{itemize}[leftmargin=1.3em]
  \item \(x\) points forward,
  \item \(y\) points left,
  \item \(z\) points upward.
\end{itemize}

The joint vector is
\[
\q =
\begin{bmatrix}
q_1 \\ q_2 \\ q_3
\end{bmatrix}
=
\begin{bmatrix}
\texttt{shoulder\_pan} \\
\texttt{shoulder\_lift} \\
\texttt{elbow\_flex}
\end{bmatrix}.
\]

\section{The intact MuJoCo model}
The MuJoCo body tree is:
\[
\begin{aligned}
\text{world} &\rightarrow \text{base} \rightarrow \text{shoulder}
\rightarrow \text{upper\_arm} \\
&\rightarrow \text{lower\_arm} \rightarrow \text{wrist}
\rightarrow \text{gripper} \rightarrow \text{moving\_jaw\_so101\_v1}.
\end{aligned}
\]

Only the first three joints are actuated:
\begin{center}
\begin{tabular}{@{}llll@{}}
\toprule
Index & Joint name & Type & Range \\
\midrule
1 & shoulder pan & revolute hinge & \([-1.91986,\;1.91986]\) rad, \(\pm110.0^\circ\) \\
2 & shoulder lift & revolute hinge & \([-1.7453293,\;1.7453293]\) rad, \(\pm100.0^\circ\) \\
3 & elbow flex & revolute hinge & \([-1.69,\;1.69]\) rad, \(\pm96.8^\circ\) \\
\bottomrule
\end{tabular}
\end{center}

The wrist and gripper bodies remain in the XML as fixed structure. The joint
count is still three, but the visual and collision geometry keeps the original
SO101 look.

The relevant fixed offsets in the MuJoCo body tree are:
\begin{center}
\begin{tabular}{@{}lll@{}}
\toprule
From & To & Translation in parent frame [m] \\
\midrule
base & shoulder & \((0.0388353,\;0,\;0.0624)\) \\
shoulder & upper\_arm & \((-0.0303992,\;-0.0182778,\;-0.0542)\) \\
upper\_arm & lower\_arm & \((-0.11257,\;-0.028,\;0)\) \\
lower\_arm & wrist & \((-0.1349,\;0.0052,\;0)\) \\
wrist & gripper & \((5.55112\text{e-}17,\;-0.0611,\;0.0181)\) \\
gripper & gripperframe site & \((0.012,\;-0.000218,\;-0.098127)\) \\
\bottomrule
\end{tabular}
\end{center}

These are exact simulator offsets, not the simplified DH parameters. They keep
the imported STL parts visually coherent.

\section{Rigid transforms}
A rigid transform is represented by a homogeneous matrix
\[
\T =
\begin{bmatrix}
\R & \p \\
0\;0\;0 & 1
\end{bmatrix},
\]
where \(\R\) is a \(3 \times 3\) rotation matrix and \(\p\) is a position vector.

For a revolute joint about the local \(z\)-axis,
\[
\Rotz(q) =
\begin{bmatrix}
\cos q & -\sin q & 0 & 0 \\
\sin q & \cos q & 0 & 0 \\
0 & 0 & 1 & 0 \\
0 & 0 & 0 & 1
\end{bmatrix}.
\]

The exact MuJoCo FK is therefore a product of fixed body transforms and three
joint rotations:
\[
\T_{0,\mathrm{ee}}(\q)
=
\T_{0,b}
\T_{b,s}^{0}\Rotz(q_1)
\T_{s,u}^{0}\Rotz(q_2)
\T_{u,l}^{0}\Rotz(q_3)
\T_{l,w}^{0}
\T_{w,g}^{0}
\T_{g,\mathrm{ee}}^{0}.
\]

In the code, FK and IK use the \texttt{gripperframe} site as the end-effector
reference point.

\section{Why a separate DH model exists}
The STL meshes and SO101 body frames were not authored as textbook DH frames.
If the XML body tree is forcibly rewritten as DH frames, the visual mesh origins
no longer line up and the robot can look disjointed.

For teaching, a separate DH model is cleaner:
\begin{itemize}[leftmargin=1.3em]
  \item the table is short,
  \item FK is a loop over rows,
  \item the Jacobian is easy to derive,
  \item IK can be taught using planar two-link geometry after the base yaw.
\end{itemize}

\section{Standard DH convention}
This project uses standard DH:
\[
\T_{i-1,i}
=
\Rotz(\theta_i)\Transz(d_i)\Transx(a_i)\Rotx(\alpha_i).
\]

The four DH symbols mean:
\begin{center}
\begin{tabular}{@{}ll@{}}
\toprule
Symbol & Meaning \\
\midrule
\(\theta_i\) & joint rotation about \(z_{i-1}\); here this is \(q_i\) \\
\(d_i\) & offset along \(z_{i-1}\) \\
\(a_i\) & link length along \(x_i\) \\
\(\alpha_i\) & link twist about \(x_i\) \\
\bottomrule
\end{tabular}
\end{center}

\section{Deriving the DH parameters}
The DH table is a teaching simplification of the intact robot. The goal is to
keep the original 3-link arm structure while making the math transparent.

The derivation is easiest to read from the joint axes:
\begin{enumerate}[leftmargin=1.3em]
  \item Joint 1 is the base yaw. Its axis is vertical, so we choose \(z_0\)
        along the shoulder pan axis. The next joint axis is orthogonal to it,
        which gives the standard DH twist \(\alpha_1 = +\pi/2\).
  \item The shoulder pitch and elbow pitch are parallel hinge axes. The
        distance between their axes along the arm centerline is \(117.5\) mm,
        so the second link length is \(a_2 = 0.1175\) m and \(\alpha_2 = 0\).
  \item The elbow pitch to fixed tool/wrist line is another parallel pair with
        a centerline separation of \(95.0\) mm, so \(a_3 = 0.0950\) m and
        \(\alpha_3 = 0\).
  \item The shoulder axis sits \(59.5\) mm above the world origin, which is
        captured by \(d_1 = 0.0595\) m. The remaining links are modeled as
        planar offsets, so \(d_2 = d_3 = 0\).
  \item The gripper/tool center lies another \(10\) mm along the final
        \(x\)-axis, which becomes the teaching tool offset \(a_{tcp}=0.0100\)
        m.
\end{enumerate}

In other words, the DH model keeps the arm's three active axes and the gross
link lengths, but it absorbs the detailed wrist and gripper shape into a fixed
tool offset.

\section{Teaching DH table}
The teaching DH model uses:
\begin{center}
\begin{tabular}{@{}llllll@{}}
\toprule
\(i\) & Joint & \(\theta_i\) & \(d_i\) [m] & \(a_i\) [m] & \(\alpha_i\) \\
\midrule
1 & shoulder pan & \(q_1\) & ? & ? & ? \\
2 & shoulder lift & \(q_2\) & ? & ? & ? \\
3 & elbow flex & \(q_3\) & ? & ? & ? \\
\bottomrule
\end{tabular}
\end{center}

The TCP/tool offset is
\[
a_{tcp} = ? \text{ m}
\]
along the final \(x\)-axis. This represents the 10 mm ball radius/tool offset in
the teaching model.

\section{DH forward kinematics}
The full DH transform is
\[
\T_{0,tcp}^{DH}(\q)
=
\T_{0,1}(q_1)\T_{1,2}(q_2)\T_{2,3}(q_3)\Transx(a_{tcp}).
\]

For the TCP position, this can be reduced to a set of trigonometric equations for \(x\), \(y\), and \(z\) in terms of the joint angles and link lengths. (Derive these as part of your exercise!)

\section{Workspace limits}
For the DH teaching model, the moving part after the base yaw is a planar
two-link arm:
\[
L_2 = a_2 = 0.1175\text{ m}, \qquad
L_3 = a_3+a_{tcp} = 0.1050\text{ m}.
\]

Ignoring joint limits, the TCP distance from the shoulder pitch frame must obey
\[
|L_2-L_3| \leq \sqrt{x^2+y^2+(z-d_1)^2} \leq L_2+L_3.
\]

Numerically:
\[
0.0125\text{ m} \leq
\sqrt{x^2+y^2+(z-0.0595)^2}
\leq 0.2225\text{ m}.
\]

So the idealized maximum reach from the shoulder pitch frame is
\[
0.2225\text{ m} = 222.5\text{ mm}.
\]

The base yaw sweeps this planar workspace around the vertical axis, but only
through its joint range:
\[
q_1 \in [-1.91986,\;1.91986]\text{ rad}
= [-110.0^\circ,\;110.0^\circ].
\]

Without joint limits, the TCP height is bounded by
\[
z_{min} = 0.0595 - 0.2225 = -0.1630\text{ m}
\]
and
\[
z_{max} = 0.0595 + 0.2225 = 0.2820\text{ m}.
\]

With the actual joint limits,
\[
\begin{aligned}
q_2 &\in [-1.7453293,\;1.7453293]\text{ rad}
= [-100.0^\circ,\;100.0^\circ],\\
q_3 &\in [-1.69,\;1.69]\text{ rad}
= [-96.8^\circ,\;96.8^\circ],
\end{aligned}
\]
the real reachable region is a subset of that ideal volume. A target can be
inside the ideal distance bound and still fail if it requires a forbidden joint
angle.

For inverse kinematics inputs, it is often useful to know approximate
axis-aligned bounds of reachable TCP positions. With the DH teaching model and
the joint limits above, a grid/optimization search gives:
\[
\begin{aligned}
x &\in [-0.1254,\;0.2225]\text{ m},\\
y &\in [-0.2225,\;0.2225]\text{ m},\\
z &\in [-0.1630,\;0.2820]\text{ m}.
\end{aligned}
\]

Equivalently, in millimeters:
\[
\begin{aligned}
x &\in [-125.4,\;222.5]\text{ mm},\\
y &\in [-222.5,\;222.5]\text{ mm},\\
z &\in [-163.0,\;282.0]\text{ mm}.
\end{aligned}
\]

These are independent coordinate bounds, not a rectangular box in which every
point is reachable. A target such as \((0.22,0.22,0.28)\) m is inside the
coordinate ranges but is still unreachable because it violates the arm's radial
reach. The negative \(x\) bound is possible because the signed planar term
\(r\) can become negative when the elbow folds back within its joint limits.

\section{Jacobian}
The position Jacobian maps joint velocity to TCP linear velocity:
\[
\dot{\p}_{tcp} = \J_p(\q)\dot{\q}.
\]

For the DH teaching model, derive the analytical Jacobian by taking the partial derivatives of your forward kinematics position equations with respect to each joint variable.

\section{Inverse kinematics}
The IK problem asks for \(\q\) given a target end-effector position
\[
\p_d = [x_d,\;y_d,\;z_d]^T.
\]

Hint: The base joint \(q_1\) can be found from the horizontal projection of the target.

Then the remaining two joints solve a planar two-link problem. You can use the law of cosines to find the two possible elbow branches for \(q_3\), and then solve for \(q_2\). (Derive the exact equations as part of your exercise!)

\section{Validation}
Current checks:
\begin{itemize}[leftmargin=1.3em]
  \item The intact MuJoCo model loads and steps with three joints and three actuators.
  \item Exact simulator FK is checked by reading the \texttt{gripperframe} site after setting joint values.
  \item The IK solver uses an analytical derivation based purely on the DH parameters.
\end{itemize}

The important interpretation is:
\begin{itemize}[leftmargin=1.3em]
  \item MuJoCo FK is exact for the simulation model.
  \item DH FK/IK is exact for the teaching model.
  \item They are intentionally separate so the visual robot stays intact.
\end{itemize}

\section{Keyframes}
The model includes:
\begin{center}
\begin{tabular}{@{}lll@{}}
\toprule
Keyframe & Joint values [rad] & Joint values [deg] \\
\midrule
home & \([0,\;0,\;0]\) & \([0.0,\;0.0,\;0.0]\) \\
ready & \([0,\;0.6,\;-0.8]\) & \([0.0,\;34.4,\;-45.8]\) \\
reach forward & \([0,\;1.0,\;-1.0]\) & \([0.0,\;57.3,\;-57.3]\) \\
reach left & \([1.2,\;0.6,\;-0.8]\) & \([68.8,\;34.4,\;-45.8]\) \\
folded & \([0,\;-1.76,\;1.74]\) & \([0.0,\;-100.8,\;99.7]\) \\
\bottomrule
\end{tabular}
\end{center}

\section{Summary}
The project should be read as:
\[
\text{intact SO101 MuJoCo model} + \text{clean DH teaching model}.
\]

This keeps the robot visually recognizable and physically simulatable while
still giving students a clear DH table, FK derivation, Jacobian, and IK solver.

\end{document}
